\documentclass{exam-cau} % 选用 Windows 字体集编译

% 1. 全局试卷元数据配置
\setyear{2025}
\setsemester{秋季学期}
\setsubject{高等数学A}
\setTotalExNumber{24} % 假设全卷共 24 道小题

% 2. 状态开关触发（根据需要取消注释）
\includeTable   % 显示顶部的得分表
\includeNotice  % 显示注意事项
\includeAnswer  % 开启此命令将同时编译出答案（发给学生时请注释掉此行）
\includeProblemCount % 开启此命令则显示大题总数 

\begin{document}

% 3. 生成试卷大标题
\generateExamTitle


% 4. 开始编写大题内容
\problem{共20分，每小题4分}{选择题}{下列各题的四个选项中，只有一个选项是正确的，请将正确答案涂在答题卡上。

  \begin{enumerate}
    \item 设函数 $f(x) = \frac{\sin x}{x}$，则 $\lim_{x \to 0} f(x) =$ \score{4}
    \begin{answer}
         解：根据重要极限可知 $\lim_{x \to 0} \frac{\sin x}{x} = 1$。故选 A。
    \end{answer}

  \item 设 $y = x^2$，则其导数 $y' =$ \score{4}
    \begin{answer}
        解：$y' = 2x$。
    \end{answer}
  \end{enumerate}
}

\problem{共30分}{计算题}{解答应写出文字说明、证明过程或演算步骤。
  \begin{enumerate}
    \item 计算定积分 $\int_a^b x^2 dx$.

    \begin{answer}
      第一步是求出原函数为 $\frac{1}{3}x^3$ \score{6}
      第二步带入牛顿-莱布尼茨公式上下限得出最终结果 \score{4}
    \end{answer}
  \end{enumerate}
}
  
% 5. 生成所有大题 
\generateProblemTable

\end{document}
